\chapter{市场状态与机会架构}

市场层面的结论是“选择性再定价”而非“全面抬估值”。行业资金状态只能回答资本是否开始关注，不能替代公司盈利质量、现金流或估值分母。我们将行业状态与公司级证据分开处理，以避免把短期资金流误解为可配置的基本面信号。

\noindent
\begin{exhibitbox}[图表 4｜31个一级行业的阶段划分与投资含义]
\small
\begin{tabularx}{\textwidth}{L{1.9cm}L{3.5cm}X L{4.2cm}}
\toprule
\textbf{阶段} & \textbf{观察到的市场特征} & \textbf{代表行业} & \textbf{投资含义} \\
\midrule
静默吸纳 & 日/周资金同向流入且价格、估值仍受约束 & 农林牧渔、房地产 & 优先寻找利润质量与现金流同步改善的公司 \\
启动确认 & 资金与价格已出现确认，静态低位优势下降 & 石油石化、煤炭 & 关注利润兑现，不追逐单日放量 \\
资金观察 & 周度资金存在但价格、估值或日度流向尚未闭环 & 计算机、银行、传媒、国防军工 & 只做事件验证，不根据主题直接配置 \\
单日反弹 & 单日流入未获得周度确认 & 美容护理、轻工制造、食品饮料、商贸零售、钢铁、纺织服饰、社会服务、环保、汽车、建筑装饰、医药生物、通信、建筑材料、基础化工、有色金属、电力设备、电子 & 不把单日反弹视为资本布局 \\
\bottomrule
\end{tabularx}
\normalsize
\sourcenote{申万一级行业官方指数历史与公开行业资金表。行业资金观察截至2026年7月14日，周度观察截至2026年7月10日。}
\end{exhibitbox}

\section{从行业阶段到公司动作}

静默吸纳并不自动意味着低风险；例如房地产的行业PB处于低分位，但公司层面仍须通过结算、税费、债务和现金流检验。启动确认也不自动意味着高估；石油石化和煤炭的价格确认必须与价差、库存、分红和现金流相匹配。因此，公司动作由三重门槛决定：行业阶段、盈利可持续性、当前价格相对情景价值。

\noindent
\begin{tabularx}{\textwidth}{L{3.2cm}X L{4.0cm}}
\toprule
\textbf{问题} & \textbf{判断规则} & \textbf{若不成立} \\
\midrule
行业是否值得跟踪？ & 资金阶段、估值分位和价格位置至少有两项支持。 & 只保留为市场观察，不进入公司验证。 \\
利润是否可进入估值？ & 扣非为正、现金流不明显背离、无低基数或结算型扭曲。 & 启动恢复模型或保留为边界标的。 \\
价格是否可承受？ & 概率值相对现价有安全边际，且外部锚/分母时效可复核。 & 降级为事件验证或估值纪律观察。 \\
\bottomrule
\end{tabularx}
